\chapter{Platonic Symmetry: Admissible Form and Nonlinear Selection}

\section{Projectors as a Language of Form}

The sphere provides a setting in which the concept of invariant form can be made exact without pretending that the geometry alone supplies the dynamics. Let

\[
\mathcal H=L^2(S^2)
\]

or, when complex amplitudes are required, $\mathcal H_{\mathbb C}=L^2(S^2,\mathbb C)$. A rotation $g$ acts through

\[
(U_g\Psi)(\hat n)=\Psi(g^{-1}\hat n).
\]

For a finite rotation group $G$, the Reynolds projector is

\[
P_G=\frac1{|G|}\sum_{g\in G}U_g.
\]

Because $P_G$ averages over the group action, it is an orthogonal projector onto the $G$-fixed subspace:

\[
P_G^2=P_G,\qquad P_G^\dagger=P_G.
\]

This allows symmetry to be measured rather than inferred from visual resemblance. Define

\[
\mathscr C_G(\Psi)=\frac{\|P_G\Psi\|^2}{\|\Psi\|^2}
\]

and

\[
\mathscr E_G(\Psi)=\frac{\|(I-P_G)\Psi\|^2}{\|\Psi\|^2}.
\]

Orthogonality yields the exact decomposition

\[
\mathscr C_G+\mathscr E_G=1.
\]

The identity is elementary but conceptually decisive. A field can be decomposed into the portion that belongs to the chosen invariant form and the portion that violates it. In this context ``Platonic form'' does not mean a mystical geometric substance. It means membership, exact or approximate, in a fixed subspace under a finite rotation group.

The same structure supplies a hierarchy under nested groups. If $H\subset G$, then every $G$-invariant state is automatically $H$-invariant, so

\[
\operatorname{Fix}(G)\subseteq\operatorname{Fix}(H).
\]

For aligned orthogonal projectors,

\[
P_HP_G=P_G
\]

and therefore

\[
\mathscr C_H\ge \mathscr C_G.
\]

A state can possess a weaker symmetry without possessing a stronger one. The mathematics again replaces a vague word such as ``symmetric'' with an ordered family of precisely defined claims.

\section{Character Multiplicity and Admissibility}

At a fixed spherical-harmonic degree $\ell$, the question of whether $G$-invariant form is possible is answered by character theory. Let $V_\ell$ be the irreducible $SO(3)$ representation of degree $\ell$. The multiplicity of the trivial representation after restriction to the finite group $G$ is

\[
d_\ell^G=\frac1{|G|}\sum_{g\in G}\chi_\ell(g),
\]

with the $SO(3)$ character for a rotation by angle $\theta$ given by

\[
\chi_\ell(\theta)=\frac{\sin[(\ell+\frac12)\theta]}{\sin(\theta/2)}.
\]

For the chiral tetrahedral, octahedral, and icosahedral groups, the first nonconstant invariant degrees are $\ell=3$, $\ell=4$, and $\ell=6$ respectively. Through $\ell=6$, the invariant multiplicities are

\begin{center}
\begin{tabular}{c c c c}
\toprule
$\ell$ & $d_\ell^T$ & $d_\ell^O$ & $d_\ell^I$ \\
\midrule
0 & 1 & 1 & 1 \\
1 & 0 & 0 & 0 \\
2 & 0 & 0 & 0 \\
3 & 1 & 0 & 0 \\
4 & 1 & 1 & 0 \\
5 & 0 & 0 & 0 \\
6 & 2 & 1 & 1 \\
\bottomrule
\end{tabular}
\end{center}

The zero entries are more powerful than the nonzero entries. If $d_\ell^G=0$, a nonzero degree-$\ell$ field cannot possess $G$ symmetry. The mathematics excludes the form. If $d_\ell^G>0$, the representation merely contains room for such a form. A two-dimensional invariant subspace, as occurs for tetrahedral symmetry at $\ell=6$, does not identify one unique ``tetrahedral shape''; it identifies a space of invariant states. Mathematical admissibility is therefore weaker than dynamical selection even before nonlinear stability is considered.

This asymmetry is one of the structural themes of the book. Formal mathematics often has extraordinary power to rule out a proposal whose required representation does not exist, but the existence of a representation is only the beginning of evidence. The same logic later governs consciousness: a theory whose mathematics cannot represent a claimed phenomenal property can fail at the level of admissibility, while a theory whose mathematics can represent it has only survived the first gate.

\section{$SO(3)$, $O(3)$, and Nonlinear Order}

The nonlinear language depends on the symmetry group. For real spherical harmonics,

\[
A(-\hat n)=(-1)^\ell A(\hat n).
\]

The symmetric square decomposes as

\[
\operatorname{Sym}^2(V_\ell)=V_0\oplus V_2\oplus\cdots\oplus V_{2\ell}.
\]

It follows that

\[
V_\ell\subset\operatorname{Sym}^2(V_\ell)
\quad\Longleftrightarrow\quad
\ell\ \text{is even}.
\]

For odd $\ell$,

\[
[\operatorname{Sym}^3(V_\ell)]^{SO(3)}=0,
\]

so there is no scalar cubic invariant under $SO(3)$. If inversion is included and the full group is $O(3)$, odd $\ell$ additionally satisfies $F(A)=F(-A)$, which removes every odd order from an invariant reduced energy. These two statements must remain distinct: $SO(3)$ eliminates the cubic scalar invariant in odd degree, while $O(3)$ imposes the stronger parity condition on all odd orders.

The distinction matters because the symmetry assumption controls the terms available to the reduced dynamics. A mathematical language does not merely record structure after the fact; its group action determines which nonlinear couplings can even appear.

\section{From Admissibility to Selection}

Consider the real rotationally equivariant equation

\[
\partial_t\Psi=\mu\Psi-\xi[\Delta_{S^2}+\ell_0(\ell_0+1)]^2\Psi+\nu\Psi^2-g\Psi^3.
\]

It is gradient flow for the energy

\[
\mathcal F[\psi]=\int_{S^2}\left[-\frac\mu2\psi^2+\frac\xi2\left([\Delta+\ell_0(\ell_0+1)]\psi\right)^2-\frac\nu3\psi^3+\frac g4\psi^4\right]d\Omega,
\]

and hence

\[
\frac{d\mathcal F}{dt}=-\int_{S^2}|\partial_t\psi|^2d\Omega\le0.
\]

For the baseline model $\mu=1/5$, $\xi=1/50$, $\nu=g=1$, repeated computations in the tested parameter regime produce an octahedral branch in the $\ell=4$ sector and an icosahedral branch in the $\ell=6$ sector. The measured positive transverse Hessian gaps are approximately

\[
\gamma_4\approx0.26207684,
\qquad
\gamma_6\approx0.55644312.
\]

These results are evidence for local stability modulo rotations in that model. They are not universal Platonic laws. A local branch must still be distinguished from a global minimum and from the probability of actualization under a distribution of initial conditions. The correct hierarchy is therefore admissibility, stationarity, local persistence, global preference, basin structure, and actualization frequency. Once stated, the hierarchy need not be repeated as a disclaimer after every calculation; it is the standing rule under which the calculations are interpreted.

The $\ell=5$ sector becomes especially informative because quartic order does not isolate a unique intrinsic shape. Instead of treating that as a failed selection, the next two chapters use it to expose a mathematically exact limit of descriptive resolution.
